%%=============================================================================
%% Voorwoord
%%=============================================================================

\chapter*{\IfLanguageName{dutch}{Woord vooraf}{Preface}}%
\label{ch:voorwoord}

%% TODO:
%% Het voorwoord is het enige deel van de bachelorproef waar je vanuit je
%% eigen standpunt (``ik-vorm'') mag schrijven. Je kan hier bv. motiveren
%% waarom jij het onderwerp wil bespreken.
%% Vergeet ook niet te bedanken wie je geholpen/gesteund/... heeft

In deze bachelorproef onderzocht ik hoe een Privileged Access Management (PAM)-tool kan bijdragen aan de beveiliging van servers binnen een bedrijfsnetwerk. Meer specifiek wilde ik nagaan hoe beheertoegang via SSH- en RDP-verbindingen beter kan worden gecontroleerd, gelogd en tijdelijk toegekend om zo beveiligingsrisico’s te verkleinen of te vermijden. Daarnaast werd onderzocht hoe moderne enterprise-organisaties PAM-oplossingen strategisch inzetten om governance, compliance en gecontroleerde toegang binnen complexe IT-omgevingen te verbeteren.

Mijn interesse in dit onderwerp ontstond vanuit mijn fascinatie voor systeem- en netwerkbeheer en de groeiende aandacht voor cybersecurity binnen moderne IT-omgevingen. Tijdens mijn opleiding werd duidelijk dat accounts met verhoogde rechten een cruciale rol spelen in zowel het beheer als de beveiliging van infrastructuur. Tegelijk vormen ze een van de grootste risico’s wanneer toegang onvoldoende gecontroleerd wordt. Het evenwicht zoeken tussen noodzakelijke toegang en beveiligingsrisico’s motiveerde mij om dit thema verder uit te werken.

Tijdens het onderzoek werd bovendien duidelijk dat Privileged Access Management binnen grote organisaties veel verder gaat dan enkel password management. De casestudies binnen deze bachelorproef tonen aan hoe enterprise-organisaties zoals Ahold Delhaize PAM-capabilities prioriteren op basis van risico, governance, operationele impact en gebruikservaring, en hoe deze strategische keuzes vervolgens vertaald worden naar praktische implementaties binnen operationele IT omgevingen.

Daarnaast wordt het belang van correcte toegangscontrole steeds groter door strengere regelgeving en de toenemende professionalisering van cyberaanvallen. Privileged Access Management biedt oplossingen voor zowel technische beveiliging enerzijds en het voldoen aan nieuwe wetgeving en regelgeving anderzijds. Ook principes zoals least privilege, Just-in-Time toegang en auditing spelen hierbij een steeds belangrijkere rol binnen moderne hybride infrastructuren.

Met deze bachelorproef hoop ik niet alleen mijn eigen kennis te verdiepen, maar ook een praktisch en bruikbaar inzicht te bieden voor professionals die hun serverinfrastructuur veiliger willen beheren en beter inzicht willen krijgen in de strategische en operationele keuzes achter moderne PAM-oplossingen.

Graag wil ik mijn promotoren Mvr. N. Declercq, Mvr. L. De Mol en co-promotor Mvr. L. Gillis bedanken voor de begeleiding, feedback en ondersteuning tijdens het uitwerken van deze bachelorproef.  Daarnaast wil ik ook Dhr. J. Hermans en Dhr. H. Populaire van Ahold Delhaize Group bedanken voor het delen van hun expertise, inzichten en praktijkervaring rond Privileged Access Management. Hun ondersteuning en de waardevolle informatie die zij ter beschikking stelden, hebben een belangrijke bijdrage geleverd aan dit onderzoek. Tot slot wil ik mijn familie en vrienden bedanken voor hun voortdurende steun en aanmoediging gedurende mijn opleiding.



